% Options for packages loaded elsewhere
\PassOptionsToPackage{unicode}{hyperref}
\PassOptionsToPackage{hyphens}{url}
%
\documentclass[
  ignorenonframetext,
]{beamer}
\usepackage{pgfpages}
\setbeamertemplate{caption}[numbered]
\setbeamertemplate{caption label separator}{: }
\setbeamercolor{caption name}{fg=normal text.fg}
\beamertemplatenavigationsymbolsempty
% Prevent slide breaks in the middle of a paragraph
\widowpenalties 1 10000
\raggedbottom
\setbeamertemplate{part page}{
  \centering
  \begin{beamercolorbox}[sep=16pt,center]{part title}
    \usebeamerfont{part title}\insertpart\par
  \end{beamercolorbox}
}
\setbeamertemplate{section page}{
  \centering
  \begin{beamercolorbox}[sep=12pt,center]{part title}
    \usebeamerfont{section title}\insertsection\par
  \end{beamercolorbox}
}
\setbeamertemplate{subsection page}{
  \centering
  \begin{beamercolorbox}[sep=8pt,center]{part title}
    \usebeamerfont{subsection title}\insertsubsection\par
  \end{beamercolorbox}
}
\AtBeginPart{
  \frame{\partpage}
}
\AtBeginSection{
  \ifbibliography
  \else
    \frame{\sectionpage}
  \fi
}
\AtBeginSubsection{
  \frame{\subsectionpage}
}
\usepackage{amsmath,amssymb}
\usepackage{lmodern}
\usepackage{iftex}
\ifPDFTeX
  \usepackage[T1]{fontenc}
  \usepackage[utf8]{inputenc}
  \usepackage{textcomp} % provide euro and other symbols
\else % if luatex or xetex
  \usepackage{unicode-math}
  \defaultfontfeatures{Scale=MatchLowercase}
  \defaultfontfeatures[\rmfamily]{Ligatures=TeX,Scale=1}
\fi
\usetheme[]{Montpellier}
% Use upquote if available, for straight quotes in verbatim environments
\IfFileExists{upquote.sty}{\usepackage{upquote}}{}
\IfFileExists{microtype.sty}{% use microtype if available
  \usepackage[]{microtype}
  \UseMicrotypeSet[protrusion]{basicmath} % disable protrusion for tt fonts
}{}
\makeatletter
\@ifundefined{KOMAClassName}{% if non-KOMA class
  \IfFileExists{parskip.sty}{%
    \usepackage{parskip}
  }{% else
    \setlength{\parindent}{0pt}
    \setlength{\parskip}{6pt plus 2pt minus 1pt}}
}{% if KOMA class
  \KOMAoptions{parskip=half}}
\makeatother
\usepackage{xcolor}
\IfFileExists{xurl.sty}{\usepackage{xurl}}{} % add URL line breaks if available
\IfFileExists{bookmark.sty}{\usepackage{bookmark}}{\usepackage{hyperref}}
\hypersetup{
  hidelinks,
  pdfcreator={LaTeX via pandoc}}
\urlstyle{same} % disable monospaced font for URLs
\newif\ifbibliography
\usepackage{color}
\usepackage{fancyvrb}
\newcommand{\VerbBar}{|}
\newcommand{\VERB}{\Verb[commandchars=\\\{\}]}
\DefineVerbatimEnvironment{Highlighting}{Verbatim}{commandchars=\\\{\}}
% Add ',fontsize=\small' for more characters per line
\usepackage{framed}
\definecolor{shadecolor}{RGB}{248,248,248}
\newenvironment{Shaded}{\begin{snugshade}}{\end{snugshade}}
\newcommand{\AlertTok}[1]{\textcolor[rgb]{0.94,0.16,0.16}{#1}}
\newcommand{\AnnotationTok}[1]{\textcolor[rgb]{0.56,0.35,0.01}{\textbf{\textit{#1}}}}
\newcommand{\AttributeTok}[1]{\textcolor[rgb]{0.77,0.63,0.00}{#1}}
\newcommand{\BaseNTok}[1]{\textcolor[rgb]{0.00,0.00,0.81}{#1}}
\newcommand{\BuiltInTok}[1]{#1}
\newcommand{\CharTok}[1]{\textcolor[rgb]{0.31,0.60,0.02}{#1}}
\newcommand{\CommentTok}[1]{\textcolor[rgb]{0.56,0.35,0.01}{\textit{#1}}}
\newcommand{\CommentVarTok}[1]{\textcolor[rgb]{0.56,0.35,0.01}{\textbf{\textit{#1}}}}
\newcommand{\ConstantTok}[1]{\textcolor[rgb]{0.00,0.00,0.00}{#1}}
\newcommand{\ControlFlowTok}[1]{\textcolor[rgb]{0.13,0.29,0.53}{\textbf{#1}}}
\newcommand{\DataTypeTok}[1]{\textcolor[rgb]{0.13,0.29,0.53}{#1}}
\newcommand{\DecValTok}[1]{\textcolor[rgb]{0.00,0.00,0.81}{#1}}
\newcommand{\DocumentationTok}[1]{\textcolor[rgb]{0.56,0.35,0.01}{\textbf{\textit{#1}}}}
\newcommand{\ErrorTok}[1]{\textcolor[rgb]{0.64,0.00,0.00}{\textbf{#1}}}
\newcommand{\ExtensionTok}[1]{#1}
\newcommand{\FloatTok}[1]{\textcolor[rgb]{0.00,0.00,0.81}{#1}}
\newcommand{\FunctionTok}[1]{\textcolor[rgb]{0.00,0.00,0.00}{#1}}
\newcommand{\ImportTok}[1]{#1}
\newcommand{\InformationTok}[1]{\textcolor[rgb]{0.56,0.35,0.01}{\textbf{\textit{#1}}}}
\newcommand{\KeywordTok}[1]{\textcolor[rgb]{0.13,0.29,0.53}{\textbf{#1}}}
\newcommand{\NormalTok}[1]{#1}
\newcommand{\OperatorTok}[1]{\textcolor[rgb]{0.81,0.36,0.00}{\textbf{#1}}}
\newcommand{\OtherTok}[1]{\textcolor[rgb]{0.56,0.35,0.01}{#1}}
\newcommand{\PreprocessorTok}[1]{\textcolor[rgb]{0.56,0.35,0.01}{\textit{#1}}}
\newcommand{\RegionMarkerTok}[1]{#1}
\newcommand{\SpecialCharTok}[1]{\textcolor[rgb]{0.00,0.00,0.00}{#1}}
\newcommand{\SpecialStringTok}[1]{\textcolor[rgb]{0.31,0.60,0.02}{#1}}
\newcommand{\StringTok}[1]{\textcolor[rgb]{0.31,0.60,0.02}{#1}}
\newcommand{\VariableTok}[1]{\textcolor[rgb]{0.00,0.00,0.00}{#1}}
\newcommand{\VerbatimStringTok}[1]{\textcolor[rgb]{0.31,0.60,0.02}{#1}}
\newcommand{\WarningTok}[1]{\textcolor[rgb]{0.56,0.35,0.01}{\textbf{\textit{#1}}}}
\setlength{\emergencystretch}{3em} % prevent overfull lines
\providecommand{\tightlist}{%
  \setlength{\itemsep}{0pt}\setlength{\parskip}{0pt}}
\setcounter{secnumdepth}{-\maxdimen} % remove section numbering
%\documentclass[unicode,10pt]{beamer}% 'unicode'が必要
\usepackage{luatexja}% 日本語を使う場合は必須
%\usepackage[japanese]{babel}	
%\usefonttheme[onlymath]{serif}
%\usepackage[T1]{fontenc} %これを使うと、ギリシャ文字が出ない
%\usepackage{textcomp}
%\usepackage[scale=1.0]{tgheros} % Sans serif font
%\usepackage[scaled]{beramono} % Monospace font
%\usepackage{luatexja-otf}
%\renewcommand{\kanjifamilydefault}{\gtdefault}
%\usepackage[haranoaji]{luatexja-preset}

% スライド番号の表示 %%%%%%%%%%%%%%%%%%%
\makeatletter
\setbeamertemplate{footline}{%
  \leavevmode%
  \hbox{%
  \begin{beamercolorbox}[wd=.5\paperwidth,ht=2.25ex,dp=1ex,right]{date in head/foot}%
    \insertframenumber{} \hspace*{2ex} % new version without total frames
  \end{beamercolorbox}}%
  \vskip0pt%
}
\makeatother
% スライド番号の表示 %%%%%%%%%%%%%%%%%%%
\ifLuaTeX
  \usepackage{selnolig}  % disable illegal ligatures
\fi

\title{第1章: イントロダクション}
\subtitle{今井耕介 著\\
『社会科学のためのデータ分析入門 (QSS)』}
\author{}
\date{\vspace{-2.5em}2026-03-09}

\begin{document}
\frame{\titlepage}

\hypertarget{rux5165ux9580}{%
\section{1.1 R入門}\label{rux5165ux9580}}

\begin{frame}[fragile]{計算操作 (Arithmetic Operations)}
\protect\hypertarget{ux8a08ux7b97ux64cdux4f5c-arithmetic-operations}{}
\begin{itemize}
\tightlist
\item
  Rは強力な計算機として使用できる。
\end{itemize}

\begin{Shaded}
\begin{Highlighting}[]
\CommentTok{\# 基本的な四則演算}
\DecValTok{5} \SpecialCharTok{+} \DecValTok{3}
\end{Highlighting}
\end{Shaded}

\begin{verbatim}
## [1] 8
\end{verbatim}

\begin{Shaded}
\begin{Highlighting}[]
\DecValTok{5} \SpecialCharTok{{-}} \DecValTok{3}
\end{Highlighting}
\end{Shaded}

\begin{verbatim}
## [1] 2
\end{verbatim}

\begin{Shaded}
\begin{Highlighting}[]
\DecValTok{5} \SpecialCharTok{/} \DecValTok{3}
\end{Highlighting}
\end{Shaded}

\begin{verbatim}
## [1] 1.666667
\end{verbatim}

\begin{Shaded}
\begin{Highlighting}[]
\CommentTok{\# べき乗と関数の使用}
\DecValTok{5} \SpecialCharTok{\^{}} \DecValTok{3}
\end{Highlighting}
\end{Shaded}

\begin{verbatim}
## [1] 125
\end{verbatim}

\begin{Shaded}
\begin{Highlighting}[]
\FunctionTok{sqrt}\NormalTok{(}\DecValTok{4}\NormalTok{)}
\end{Highlighting}
\end{Shaded}

\begin{verbatim}
## [1] 2
\end{verbatim}
\end{frame}

\hypertarget{ux30aaux30d6ux30b8ux30a7ux30afux30c8-objects}{%
\section{1.2 オブジェクト
(Objects)}\label{ux30aaux30d6ux30b8ux30a7ux30afux30c8-objects}}

\begin{frame}[fragile]{オブジェクトの作成と上書き}
\protect\hypertarget{ux30aaux30d6ux30b8ux30a7ux30afux30c8ux306eux4f5cux6210ux3068ux4e0aux66f8ux304d}{}
\begin{itemize}
\tightlist
\item
  \texttt{\textless{}-}
  を用いて、計算結果やデータを名前付きの「オブジェクト」として保存する。
\end{itemize}

\begin{Shaded}
\begin{Highlighting}[]
\NormalTok{result }\OtherTok{\textless{}{-}} \DecValTok{5} \SpecialCharTok{+} \DecValTok{3}
\NormalTok{result}
\end{Highlighting}
\end{Shaded}

\begin{verbatim}
## [1] 8
\end{verbatim}

\begin{Shaded}
\begin{Highlighting}[]
\CommentTok{\# 文字列の保存}
\NormalTok{kosuke }\OtherTok{\textless{}{-}} \StringTok{"instructor"}
\NormalTok{kosuke}
\end{Highlighting}
\end{Shaded}

\begin{verbatim}
## [1] "instructor"
\end{verbatim}

\begin{itemize}
\tightlist
\item
  オブジェクトは後から上書きできる。また、大文字と小文字は区別される（\texttt{result}
  と \texttt{Result} は別物）。
\end{itemize}
\end{frame}

\hypertarget{ux30d9ux30afux30c8ux30eb-vectors}{%
\section{1.3 ベクトル
(Vectors)}\label{ux30d9ux30afux30c8ux30eb-vectors}}

\begin{frame}[fragile]{ベクトルの作成}
\protect\hypertarget{ux30d9ux30afux30c8ux30ebux306eux4f5cux6210}{}
\begin{itemize}
\tightlist
\item
  \texttt{c()}
  関数（concatenate）を用いて、複数の要素をひとまとめにした「ベクトル」を作成する。
\end{itemize}

\begin{Shaded}
\begin{Highlighting}[]
\CommentTok{\# 世界人口データ（千人単位）}
\NormalTok{world.pop }\OtherTok{\textless{}{-}} \FunctionTok{c}\NormalTok{(}\DecValTok{2525779}\NormalTok{, }\DecValTok{3026003}\NormalTok{, }\DecValTok{3691173}\NormalTok{, }\DecValTok{4449049}\NormalTok{, }
               \DecValTok{5320817}\NormalTok{, }\DecValTok{6127700}\NormalTok{, }\DecValTok{6916183}\NormalTok{)}
\NormalTok{world.pop}
\end{Highlighting}
\end{Shaded}

\begin{verbatim}
## [1] 2525779 3026003 3691173 4449049 5320817 6127700 6916183
\end{verbatim}
\end{frame}

\begin{frame}[fragile]{要素の抽出}
\protect\hypertarget{ux8981ux7d20ux306eux62bdux51fa}{}
\begin{itemize}
\tightlist
\item
  \texttt{{[}{]}} を用いてベクトルの一部を取り出すことができる。
\end{itemize}

\begin{Shaded}
\begin{Highlighting}[]
\CommentTok{\# 2番目の要素を取り出す}
\NormalTok{world.pop[}\DecValTok{2}\NormalTok{]}
\end{Highlighting}
\end{Shaded}

\begin{verbatim}
## [1] 3026003
\end{verbatim}

\begin{Shaded}
\begin{Highlighting}[]
\CommentTok{\# 2番目と4番目を取り出す}
\NormalTok{world.pop[}\FunctionTok{c}\NormalTok{(}\DecValTok{2}\NormalTok{, }\DecValTok{4}\NormalTok{)]}
\end{Highlighting}
\end{Shaded}

\begin{verbatim}
## [1] 3026003 4449049
\end{verbatim}

\begin{Shaded}
\begin{Highlighting}[]
\CommentTok{\# 3番目以外のすべての要素を取り出す}
\NormalTok{world.pop[}\SpecialCharTok{{-}}\DecValTok{3}\NormalTok{]}
\end{Highlighting}
\end{Shaded}

\begin{verbatim}
## [1] 2525779 3026003 4449049 5320817 6127700 6916183
\end{verbatim}
\end{frame}

\hypertarget{ux95a2ux6570-functions}{%
\section{1.4 関数 (Functions)}\label{ux95a2ux6570-functions}}

\begin{frame}[fragile]{組み込み関数の使用}
\protect\hypertarget{ux7d44ux307fux8fbcux307fux95a2ux6570ux306eux4f7fux7528}{}
\begin{itemize}
\tightlist
\item
  Rには多くの関数が用意されている。
\end{itemize}

\begin{Shaded}
\begin{Highlighting}[]
\CommentTok{\# 長さ、最小値、最大値、平均}
\FunctionTok{length}\NormalTok{(world.pop)  }
\end{Highlighting}
\end{Shaded}

\begin{verbatim}
## [1] 7
\end{verbatim}

\begin{Shaded}
\begin{Highlighting}[]
\FunctionTok{min}\NormalTok{(world.pop)     }
\end{Highlighting}
\end{Shaded}

\begin{verbatim}
## [1] 2525779
\end{verbatim}

\begin{Shaded}
\begin{Highlighting}[]
\FunctionTok{max}\NormalTok{(world.pop)     }
\end{Highlighting}
\end{Shaded}

\begin{verbatim}
## [1] 6916183
\end{verbatim}

\begin{Shaded}
\begin{Highlighting}[]
\FunctionTok{mean}\NormalTok{(world.pop)    }
\end{Highlighting}
\end{Shaded}

\begin{verbatim}
## [1] 4579529
\end{verbatim}
\end{frame}

\begin{frame}[fragile]{カスタム関数の作成}
\protect\hypertarget{ux30abux30b9ux30bfux30e0ux95a2ux6570ux306eux4f5cux6210}{}
\begin{itemize}
\tightlist
\item
  自分で独自の関数を定義することもできる。
\end{itemize}

\begin{Shaded}
\begin{Highlighting}[]
\CommentTok{\# 独自の要約関数}
\NormalTok{my.summary }\OtherTok{\textless{}{-}} \ControlFlowTok{function}\NormalTok{(x)\{ }
\NormalTok{  s.out }\OtherTok{\textless{}{-}} \FunctionTok{sum}\NormalTok{(x)}
\NormalTok{  l.out }\OtherTok{\textless{}{-}} \FunctionTok{length}\NormalTok{(x)}
\NormalTok{  m.out }\OtherTok{\textless{}{-}}\NormalTok{ s.out }\SpecialCharTok{/}\NormalTok{ l.out}
\NormalTok{  out }\OtherTok{\textless{}{-}} \FunctionTok{c}\NormalTok{(s.out, l.out, m.out)}
  \FunctionTok{names}\NormalTok{(out) }\OtherTok{\textless{}{-}} \FunctionTok{c}\NormalTok{(}\StringTok{"sum"}\NormalTok{, }\StringTok{"length"}\NormalTok{, }\StringTok{"mean"}\NormalTok{)}
  \FunctionTok{return}\NormalTok{(out)}
\NormalTok{\}}
\FunctionTok{my.summary}\NormalTok{(world.pop)}
\end{Highlighting}
\end{Shaded}

\begin{verbatim}
##      sum   length     mean 
## 32056704        7  4579529
\end{verbatim}
\end{frame}

\hypertarget{ux30c7ux30fcux30bfux30d5ux30a1ux30a4ux30eb-data-files}{%
\section{1.5 データファイル (Data
Files)}\label{ux30c7ux30fcux30bfux30d5ux30a1ux30a4ux30eb-data-files}}

\begin{frame}[fragile]{CSVファイルの読み込み}
\protect\hypertarget{csvux30d5ux30a1ux30a4ux30ebux306eux8aadux307fux8fbcux307f}{}
\begin{itemize}
\tightlist
\item
  \texttt{read.csv()} 関数を用いて、外部のデータセットをRに読み込む。
\end{itemize}

\begin{Shaded}
\begin{Highlighting}[]
\CommentTok{\# UNpopデータの読み込み}
\NormalTok{UNpop }\OtherTok{\textless{}{-}} \FunctionTok{read.csv}\NormalTok{(}\StringTok{"UNpop.csv"}\NormalTok{) }

\CommentTok{\# データフレームの確認}
\FunctionTok{dim}\NormalTok{(UNpop)}
\end{Highlighting}
\end{Shaded}

\begin{verbatim}
## [1] 7 2
\end{verbatim}

\begin{Shaded}
\begin{Highlighting}[]
\FunctionTok{head}\NormalTok{(UNpop, }\AttributeTok{n =} \DecValTok{3}\NormalTok{)}
\end{Highlighting}
\end{Shaded}

\begin{verbatim}
##   year world.pop
## 1 1950   2525779
## 2 1960   3026003
## 3 1970   3691173
\end{verbatim}
\end{frame}

\begin{frame}[fragile]{データフレームからの変数の抽出}
\protect\hypertarget{ux30c7ux30fcux30bfux30d5ux30ecux30fcux30e0ux304bux3089ux306eux5909ux6570ux306eux62bdux51fa}{}
\begin{itemize}
\tightlist
\item
  \texttt{\$} 演算子を用いて、特定の列（変数）にアクセスする。
\end{itemize}

\begin{Shaded}
\begin{Highlighting}[]
\CommentTok{\# world.pop変数の抽出}
\NormalTok{UNpop}\SpecialCharTok{$}\NormalTok{world.pop}
\end{Highlighting}
\end{Shaded}

\begin{verbatim}
## [1] 2525779 3026003 3691173 4449049 5320817 6127700 6916183
\end{verbatim}

\begin{itemize}
\tightlist
\item
  特定の行や列にインデックス \texttt{{[}行,\ 列{]}}
  でアクセスすることも可能。
\end{itemize}

\begin{Shaded}
\begin{Highlighting}[]
\CommentTok{\# 1〜3行目のすべての列}
\NormalTok{UNpop[}\DecValTok{1}\SpecialCharTok{:}\DecValTok{3}\NormalTok{, ]}
\end{Highlighting}
\end{Shaded}

\begin{verbatim}
##   year world.pop
## 1 1950   2525779
## 2 1960   3026003
## 3 1970   3691173
\end{verbatim}
\end{frame}

\hypertarget{ux30d1ux30c3ux30b1ux30fcux30b8-packages}{%
\section{1.6 パッケージ
(Packages)}\label{ux30d1ux30c3ux30b1ux30fcux30b8-packages}}

\begin{frame}[fragile]{Rパッケージの利用}
\protect\hypertarget{rux30d1ux30c3ux30b1ux30fcux30b8ux306eux5229ux7528}{}
\begin{itemize}
\tightlist
\item
  \textbf{パッケージ}をインストールすることで、Rの機能を拡張できる。
\item
  例：\texttt{foreign} パッケージを用いて、Stata (\texttt{.dta}) や SPSS
  (\texttt{.sav}) などの他ソフトウェアのデータを読み書きする。
\end{itemize}

\begin{Shaded}
\begin{Highlighting}[]
\CommentTok{\# インストール（初回のみ）}
\FunctionTok{install.packages}\NormalTok{(}\StringTok{"foreign"}\NormalTok{)}

\CommentTok{\# パッケージの読み込み}
\FunctionTok{library}\NormalTok{(foreign)}
\end{Highlighting}
\end{Shaded}
\end{frame}

\hypertarget{ux307eux3068ux3081}{%
\section{1.7 まとめ}\label{ux307eux3068ux3081}}

\begin{frame}[fragile]{第1章のまとめ}
\protect\hypertarget{ux7b2c1ux7ae0ux306eux307eux3068ux3081}{}
\begin{itemize}
\tightlist
\item
  \textbf{RとRStudio}の基本的な操作方法を習得した。
\item
  \textbf{オブジェクト}と\textbf{ベクトル}を用いたデータの保存と操作方法を学んだ。
\item
  組み込み関数や自作関数を用いて、データを要約する方法を理解した。
\item
  \texttt{read.csv()}
  を用いてデータフレームを読み込み、操作する方法を学んだ。
\item
  データ分析に必要な基礎的なコーディングスキルを身につけた。
\end{itemize}
\end{frame}

\end{document}
